% eigenq research paper preamble — shared across all reports/<project>.qmd
%
% Required TeX packages (install via tlmgr if rendering locally):
%   amsmath, amssymb, amsthm, mathtools, bm, booktabs, microtype,
%   biblatex, csquotes, hyperref, cleveref, enumitem, siunitx, minted

\usepackage{amsmath}
\usepackage{amssymb}
\usepackage{mathtools}
\usepackage{bm}
\usepackage{booktabs}
\usepackage{microtype}
\usepackage{enumitem}
\usepackage{siunitx}
\usepackage{cleveref}

% Number theorem-likes within sections, share a counter so numbering is contiguous.
% Quarto's native theorem callouts use these counters automatically.
\numberwithin{equation}{section}
\numberwithin{figure}{section}
\numberwithin{table}{section}

% Probability / finance shorthand.
\newcommand{\E}{\mathbb{E}}
\newcommand{\Prob}{\mathbb{P}}
\newcommand{\Q}{\mathbb{Q}}
\newcommand{\R}{\mathbb{R}}
\newcommand{\N}{\mathbb{N}}
\newcommand{\Z}{\mathbb{Z}}
\newcommand{\Filt}{\mathcal{F}}
\newcommand{\indep}{\perp\!\!\!\perp}
\newcommand{\indic}[1]{\mathbf{1}_{\{#1\}}}
\newcommand{\eqdist}{\stackrel{d}{=}}
\newcommand{\eqas}{\stackrel{\text{a.s.}}{=}}
\renewcommand{\Pr}{\mathbb{P}}

% Visual style for the standardized "Formalization" callout — referenced from
% template.qmd as `::: {.callout-note title="Formalization"}`.
\usepackage{tcolorbox}
\tcbuselibrary{breakable, skins}

% Title block styling consistent with AQR / JFE working-paper aesthetic.
\setkomafont{title}{\bfseries\Large}
\setkomafont{author}{\normalsize}
\setkomafont{date}{\small\itshape}
\setkomafont{section}{\bfseries\large}
\setkomafont{subsection}{\bfseries\normalsize}
